\par\noindent\begin{minipage}{\columnwidth}
\centering\textbf{Algorithm 1: Proposed Hybrid Spatiotemporal Forecasting Framework}\par\smallskip
\centering\footnotesize
\begin{tabularx}{\columnwidth}{@{}lX@{}}
\toprule
\multicolumn{2}{l}{\textbf{Inputs}}\\
I. & NSSP and contemporary/historical ILINet surveillance panels.\\
II. & Target and auxiliary history available through the completed week.\\
\midrule
\multicolumn{2}{l}{\textbf{Procedure}}\\
III. & Load data; harmonize types and source-week identifiers.\\
IV. & Order each series; identify missing or invalid source values.\\
V. & Define whole-week fit, validation, and final-test boundaries.\\
VI. & Build target lags, rolling summaries, annual phase, and lag-1 auxiliary inputs. Exclude rows missing target/required features; apply week partitions; fit scaling on the fit segment.\\
VII. & Fit K-means on that segment; derive and append one-hot cluster indicators.\\
VIII. & Fit RF, GBR, and MLP; predict the later validation segment.\\
IX. & Select classical orders by validation; compute weights with (\ref{eq:weights}); freeze both. Refit the scaler, K-means, learners, and classical models on all pre-test rows; combine learner forecasts with (\ref{eq:hybrid}).\\
X. & Score all methods on matched units: HSA-weeks for primary NSSP, reporting rows for ILINet. Retain NSSP record-weighted comparisons.\\
XI. & Summarize temporal errors and descriptive state maps.\\
XII. & Repeat independently for each surveillance experiment.\\
\midrule
\multicolumn{2}{l}{\textbf{Outputs}}\\
XIII. & One-week forecasts, metrics, ablations/sensitivities, and test-period geographic summaries.\\
\bottomrule
\end{tabularx}
\end{minipage}
\par\smallskip
